\section{A pre-registered plan for Thailand}\label{sec:plan}

The evidence so far establishes the rail, the weak macroeconomic motive, the regional threat and the limits of
aggregate data. Whether Thai stablecoin flows actually carry fraud proceeds, and how much, is an empirical question
that requires data held by exchanges, the Securities and Exchange Commission and the Anti-Money Laundering Office. To
prevent the specification search that such sensitive questions invite, \cref{tab:prereg} commits in advance to six
hypotheses, the data each needs, the estimator, the primary statistic and the decision rule.

\begin{table}[!htbp]
\centering\footnotesize
\caption{Pre-registered hypotheses, data, estimators and decision rules}
\label{tab:prereg}
\begin{tabularx}{\textwidth}{L{2.5cm}L{4.35cm}L{4.1cm}X}
\toprule
Hypothesis & Data (tier) & Estimator & Decision rule \\
\midrule
H1 Excess stablecoin demand & Exchange volumes by quote asset across 20+ countries; inflation, FX volatility,
remittances, migrant stock, capital-account openness, derivatives access (1--2) & Panel regression of the stablecoin
share of volume on fundamentals; Thailand's standardised residual & Residual positive and in the top decile of a
leave-one-country-out distribution \\
\addlinespace
H2 USDT exits more than other assets & Daily purchases and on-chain withdrawals by asset on Thai platforms (2--3) &
Difference between USDT and BTC withdrawal-to-purchase ratios, controlling for volatility and offshore funding rates &
One-sided test of a positive difference at 5\%; reported net of exchange-labelled destinations \\
\addlinespace
H3 Fast pass-through concentrated in young, fan-in accounts & Account-level deposit--purchase--withdrawal cycles (3) &
Gaussian mixture on log dwell; fast-component share by account age and fan-in & Difference positive with a bootstrap
interval excluding zero, after removing withdrawals to labelled exchanges \\
\addlinespace
H4 Structuring below 50,000 baht & Deposit amounts by funded asset, ideally before and after the verification rule (3)
& Difference in responding shares (\cref{prop:bunching}); difference-in-difference if pre-rule data exist
(\cref{cor:didib}) & One-sided $p<0.05$ at 50,000 baht and no significant difference at placebo thresholds of 30,000 and
70,000 baht \\
\addlinespace
H5 Crackdowns change the composition of flows & Weekly composition series (share of fast structured deposits,
peer-to-peer premium, network mix) for Thailand and donor countries (2--3) & Synthetic control and synthetic
difference-in-differences; triple difference USDT versus BTC & Rank-based placebo $p\le0.05$ with pre-period fit no
worse than the median donor's \\
\addlinespace
H6 Thai withdrawals are more exposed than peers' & Tron transfers from labelled Thai and peer-country exchange hot
wallets; issuer freezes; OFAC list (1--3) & Directed, value-weighted three-hop exposure (\cref{eq:exposure}) & Thai
exposure above the 75th percentile of peer exchanges with a bootstrap interval excluding the peer median \\
\bottomrule
\end{tabularx}
\par\vspace{3pt}{\footnotesize\color{muted}Data tiers: 1 public; 2 commercial or academic licence; 3 partnership under the
PDPA. H3--H6 form the primary family and are adjusted with the Holm--Bonferroni procedure; H1--H2 are secondary.\par}
\end{table}

\paragraph{Interpretation}
The hypotheses are ordered to separate structure from crime. If H1 and H2 hold but H3--H6 do not, Thai stablecoin
behaviour is unusual because of market structure: the absence of domestic derivatives, the baht-quoted market and
cross-border demand. If H3--H6 also hold, the residual pattern carries the fingerprints of pass-through laundering,
and \cref{prop:bunching} converts the bunching estimate into a lower bound on structured volume, since unstructured
illicit deposits are not counted.

\paragraph{Data access and privacy}
Tier-3 analyses should follow a ``code goes to the data'' model: the replication package runs on the simulator's
synthetic ledger, which has the same schema as the exchange data required, so code can be reviewed, tested and
approved before it touches personal data inside a regulator's or exchange's secure environment. Only aggregate,
disclosure-controlled outputs leave that environment. No address of a person who has not been charged is published.

\section{Policy implications}\label{sec:policy}

Six implications follow, each conditional on the evidence that the plan would produce.

\begin{enumerate}[leftmargin=1.4em,label=\arabic*.]
\item \textbf{Monitor differences, not levels.} Supervisors and exchanges should track structuring and pass-through as
contrasts between the stablecoin channel and other channels on the same platform, with normalisations that do not
depend on the shape of the amount distribution. Levels confound crime with hassle avoidance and market structure.
\item \textbf{Make exposure screening direction-aware.} Screening that counts undirected hops to flagged addresses
mechanically flags the exchange hubs to which ordinary investors withdraw. Exposure should follow funds forward,
weight by value and decay with distance.
\item \textbf{Protect honest peer-to-peer dealers.} Honest dealers are the population most easily mistaken for mules and
the usual victims of ``triangle'' fraud, in which a criminal has a victim pay the dealer directly. Requiring that baht
payments for peer-to-peer trades come from a bank account in the buyer's verified name, and creating a fast channel to
review frozen dealer accounts, would reduce both the fraud and its collateral damage.
\item \textbf{Evaluate crackdowns with composition data and honest inference.} Chain-wide aggregates cannot show whether
a crackdown worked (\cref{sec:eventstudy}). Publishing weekly, anonymised composition indicators---the share of
deposits just below verification thresholds, the share of fast exits, the peer-to-peer premium---would make evaluation
possible without disclosing personal data.
\item \textbf{Anticipate displacement.} Fixed, publicly known thresholds create the notch that structuring exploits, and
pressure on one channel moves flows to others (\cref{prop:displacement}). Risk-based verification and coordinated action
across banks, licensed exchanges and peer-to-peer markets limit the gain from switching.
\item \textbf{Use the issuer and the ledger.} Because the region's main dollar rail has a single issuer able to freeze
tokens and a public ledger, regional cooperation with the issuer and with neighbouring authorities can act after the
conversion boundary, where Thai banks no longer can.
\end{enumerate}

\section{Limitations and ethics}\label{sec:limits}

The descriptive evidence covers the chains in the Coin Metrics community files and measures turnover by transfer
counts, not values. Address-based statistics are affected by custodial omnibus wallets. Event dates and institutional
details were compiled without access to primary legal databases and must be verified. The simulation's absolute
results depend on stylised parameters; its purpose is to rank methods and reveal failure modes, not to estimate
prevalence. Statistical anomalies are not evidence of wrongdoing by any person, dealer or exchange, and the paper names
entities only where public authorities have acted. The laundering typology is described at the level of published
FATF and UNODC material and omits operational detail.

\section{Conclusion}\label{sec:conclusion}

Thailand's stablecoin question is easy to pose and hard to answer. Public data establish that the region's dominant
digital dollar, USDT on Tron, is a high-turnover retail rail on which sanctioned activity concentrates; that Thailand,
unlike its neighbours, has little macroeconomic reason to use it heavily; and that chain-wide data cannot tell whether
crackdowns work. The answer to whether Thai rails carry scam proceeds lies in the composition of flows at the boundary
where baht become tokens. This paper supplies the estimators for that boundary---a slope-invariant difference in
bunching at the face-verification threshold, dwell-time mixtures, direction-aware exposure and synthetic control---shows
in simulation which of them work and why the common heuristics fail, and commits in advance to how the answer will be
read once the data are in hand.
